
%===============================================================================
% APPENDIX
% COMPLETE DERIVATION, ECONOMIC INTERPRETATION, AND STATISTICAL GUIDE
% BRAZIL 1964--1980 DSGE MODEL
%
% Paste this file after \appendix in the main paper, or keep the \appendix
% command below if this is the first appendix.
%
% This paste-ready version requires only amsmath and amssymb in the main preamble.
% Add BEFORE \begin{document} if your main file does not already contain them:
% \usepackage{amsmath}
% \usepackage{amssymb}
%===============================================================================

\appendix
\section{Complete Derivation and Guide to the Brazil 1964--1980 DSGE Model}
\label{app:model_bible}

\subsection{Purpose of this Appendix}
\label{app:purpose}

This appendix is the complete mathematical and economic guide to the baseline
Brazilian model used in the paper. Its purpose is to make explicit every step
that is compressed in the main text and in the Dynare implementation. In
particular, it distinguishes carefully among:

\begin{enumerate}
    \item equations that follow from private optimization;
    \item accounting identities;
    \item market-clearing conditions;
    \item policy rules;
    \item stochastic processes; and
    \item semi-structural linear approximations introduced deliberately for
    empirical tractability.
\end{enumerate}

This distinction is essential. Not every equation in a DSGE model is an
optimality condition. The household Euler equation, for example, follows from a
first-order condition. The CPI relation is an identity. The Taylor-type
interest-rate rule is a policy specification. The country-risk equation is a
reduced-form pricing wedge. Treating these objects as though they had the same
logical status would obscure both the economics and the econometrics.

The appendix reproduces \emph{exactly} the 26-equation linear system used in the
successfully running baseline Dynare file. No household foreign bond is added.
No output-gap term is added to monetary policy. No direct oil shock is added to
the Phillips curve. Taxes do not respond endogenously to debt. Public
investment remains separate from ordinary government consumption. External
debt is consolidated public-sector foreign debt.

%-------------------------------------------------------------------------------
\subsection{Notation, Units, and Timing}
\label{app:notation}
%-------------------------------------------------------------------------------

\subsubsection{Levels, log deviations, and level deviations}

Upper-case letters denote economic variables in levels. For a strictly positive
variable \(X_t\) with deterministic steady state \(\bar X>0\), define its
log-deviation by

\begin{equation}
\widehat{x}_t
\equiv
\log\left(\frac{X_t}{\bar X}\right)
\simeq
\frac{X_t-\bar X}{\bar X}.
\label{eq:hatdefinition}
\end{equation}

In the Dynare code, hats are suppressed. Thus \(c_t\), \(y_t\), \(m_t\),
\(s_t\), and so forth should be read as log deviations unless stated
otherwise.

Five variables are deliberately \emph{not} treated as ordinary log deviations:

\begin{align}
tb_t &:\quad \text{trade-balance deviation as a fraction of steady-state GDP},\\
d_t  &:\quad \text{external public-debt/GDP deviation},\\
b_t  &:\quad \text{domestic public-debt/GDP deviation},\\
sg_t &:\quad \text{seigniorage/GDP deviation},\\
\tau_t &:\quad \text{tax-revenue/GDP deviation}.
\end{align}

The reason is economic as well as mathematical. In particular, the steady-state
trade balance can be zero, so \(\log(TB_t/\overline{TB})\) would be undefined.
Consequently, \(tb_t\) is linearized in levels relative to steady-state output.

Inflation rates and interest rates,

\[
\pi_t,\quad \pi_{H,t},\quad \pi_{F,t}^{*},\quad
i_t,\quad i_t^{*},\quad rp_t,
\]

are quarterly \emph{rate deviations in decimal units}. For these variables,
\(0.01\) means one percentage point at quarterly frequency.

\subsubsection{Expectations and timing}

The operator \(\mathbb E_t[\cdot]\) denotes the mathematical expectation
conditional on information available at date \(t\).

The Dynare notation maps into the equations as

\[
x(+1)\equiv \mathbb E_t x_{t+1},
\qquad
x(-1)\equiv x_{t-1}.
\]

The model contains both predetermined and forward-looking variables. Debt
stocks, lagged inflation, lagged interest rates, and persistent exogenous
processes carry information from the past. Consumption, domestic inflation, and
the exchange rate contain forward-looking components.

%-------------------------------------------------------------------------------
\subsection{A Log-Linearization Toolkit}
\label{app:linearization_toolkit}
%-------------------------------------------------------------------------------

Before deriving the model, it is useful to collect the approximations used
repeatedly.

\subsubsection{First-order Taylor approximation}

For a differentiable function \(F(X)\),

\begin{equation}
F(X_t)
\simeq
F(\bar X)
+
F_X(\bar X)(X_t-\bar X).
\end{equation}

For a positive variable represented as \(X_t=\bar X e^{x_t}\),

\begin{equation}
X_t
\simeq
\bar X(1+x_t).
\end{equation}

\subsubsection{Products}

If \(Z_t=X_tY_t\), then to first order

\begin{equation}
z_t=x_t+y_t.
\end{equation}

Cross-products of deviations such as \(x_ty_t\) are second-order terms and are
discarded in a first-order solution.

\subsubsection{Ratios}

If \(Z_t=X_t/Y_t\), then

\begin{equation}
z_t=x_t-y_t.
\end{equation}

\subsubsection{Weighted sums}

Suppose

\[
Y_t=X_{1,t}+X_{2,t}+\cdots+X_{J,t}
\]

and let

\[
\omega_j=\frac{\bar X_j}{\bar Y}.
\]

Then the first-order approximation is

\begin{equation}
y_t
=
\sum_{j=1}^{J}\omega_j x_{j,t}.
\label{eq:sharelinearization}
\end{equation}

This rule generates the goods-market-clearing equation and the trade-balance
equation used in Dynare.

\subsubsection{Why constants disappear}

A linear model is written in deviations from a deterministic steady state.
Therefore steady-state constants vanish after subtracting the steady-state
equation. A linear debt equation can consequently contain a nonzero benchmark
debt ratio \(\bar d\) even though the Dynare variable \(d_t\) itself has
steady-state value zero: \(d_t\) is the \emph{deviation} around \(\bar d\).

%===============================================================================
\subsection{Households}
\label{app:households}
%===============================================================================

\subsubsection{Preferences}

The representative household has expected lifetime utility

\begin{equation}
\mathcal U_0
=
\mathbb E_0
\sum_{t=0}^{\infty}
\beta^t
\left[
\frac{C_t^{1-\sigma}}{1-\sigma}
-
\chi\frac{N_t^{1+\varphi}}{1+\varphi}
+
\zeta_m\frac{(M_t/P_t)^{1-\nu}}{1-\nu}
\right],
\label{eq:utility_level}
\end{equation}

where

\begin{itemize}
    \item \(0<\beta<1\) is the discount factor;
    \item \(\sigma>0\) is the coefficient of relative risk aversion and the
    inverse intertemporal elasticity of substitution under CRRA utility;
    \item \(\varphi>0\) is the inverse Frisch elasticity parameter;
    \item \(M_t/P_t\) denotes real money balances;
    \item \(\nu>0\) governs the curvature of liquidity services.
\end{itemize}

Money enters utility solely to provide a coherent demand for real balances.
It is not used to impose an ad hoc relation in which inflation directly causes
consumption.

\subsubsection{Household nominal budget constraint}

Households consume, hold money, and buy one-period domestic public debt:

\begin{equation}
P_tC_t+M_t+Q_t^dB_{t+1}
=
W_tN_t+\Pi_t-T_t^h+TR_t+M_{t-1}+B_t.
\label{eq:hh_budget_level}
\end{equation}

The object \(B_{t+1}\) is a one-period nominal bond purchased at price
\(Q_t^d\) and paying one unit of domestic currency at \(t+1\). Hence

\begin{equation}
Q_t^d=\frac{1}{1+i_t^{\mathrm{lev}}},
\label{eq:bond_price}
\end{equation}

where \(i_t^{\mathrm{lev}}\) denotes the net nominal rate in levels.

There is no household foreign bond. This is a defining feature of the final
model.

\subsubsection{The household Lagrangian}

Let \(\Lambda_t\) be the multiplier on the nominal budget constraint. The
Lagrangian is

\begin{align}
\mathcal L
=
\mathbb E_0
\sum_{t=0}^{\infty}
\beta^t
\Bigg\{&
\frac{C_t^{1-\sigma}}{1-\sigma}
-
\chi\frac{N_t^{1+\varphi}}{1+\varphi}
+
\zeta_m\frac{(M_t/P_t)^{1-\nu}}{1-\nu}
\nonumber\\
&+
\Lambda_t
\left[
W_tN_t+\Pi_t-T_t^h+TR_t+M_{t-1}+B_t
-P_tC_t-M_t-Q_t^dB_{t+1}
\right]
\Bigg\}.
\label{eq:hh_lagrangian}
\end{align}

\subsubsection{Consumption first-order condition}

Differentiating with respect to \(C_t\),

\begin{equation}
C_t^{-\sigma}
=
\Lambda_tP_t,
\label{eq:foc_c}
\end{equation}

or

\begin{equation}
\Lambda_t
=
\frac{C_t^{-\sigma}}{P_t}.
\label{eq:lambda}
\end{equation}

Economically, \(\Lambda_t\) is the marginal utility value of one additional
unit of nominal wealth.

\subsubsection{Labor-supply first-order condition}

Differentiating with respect to \(N_t\),

\begin{equation}
\chi N_t^{\varphi}
=
\Lambda_tW_t.
\end{equation}

Using equation \eqref{eq:lambda},

\begin{equation}
\frac{W_t}{P_t}
=
\chi N_t^\varphi C_t^\sigma.
\label{eq:labor_supply_level}
\end{equation}

Taking logs and removing constants gives

\begin{equation}
\boxed{
w_t-p_t
=
\sigma c_t+\varphi n_t.
}
\label{eq:labor_supply_linear}
\end{equation}

This condition is not entered as a separate Dynare equation because it is
substituted directly into the marginal-cost equation. That is why the
26-equation system contains no stand-alone wage equation.

\subsubsection{Domestic-bond first-order condition}

Differentiating with respect to \(B_{t+1}\) gives

\begin{equation}
\Lambda_tQ_t^d
=
\beta\mathbb E_t\Lambda_{t+1}.
\end{equation}

Substituting equation \eqref{eq:lambda},

\begin{equation}
Q_t^d
=
\beta
\mathbb E_t
\left[
\left(\frac{C_{t+1}}{C_t}\right)^{-\sigma}
\frac{P_t}{P_{t+1}}
\right].
\label{eq:bond_euler_price}
\end{equation}

Since \(Q_t^d=1/(1+i_t^{\mathrm{lev}})\),

\begin{equation}
1
=
\beta(1+i_t^{\mathrm{lev}})
\mathbb E_t
\left[
\left(\frac{C_{t+1}}{C_t}\right)^{-\sigma}
\frac{1}{\Pi_{t+1}}
\right],
\label{eq:euler_exact}
\end{equation}

where

\[
\Pi_{t+1}\equiv\frac{P_{t+1}}{P_t}.
\]

\subsubsection{Log-linear Euler equation}

Take a first-order approximation of equation \eqref{eq:euler_exact} around the
steady state. Constants cancel by the steady-state Euler condition. The
result is

\begin{equation}
\sigma
\left(
c_t-\mathbb E_tc_{t+1}
\right)
=
i_t-\mathbb E_t\pi_{t+1},
\end{equation}

or

\begin{equation}
\boxed{
c_t
=
\mathbb E_tc_{t+1}
-
\frac{1}{\sigma}
\left(
i_t-\mathbb E_t\pi_{t+1}
\right).
}
\label{eq:euler_final}
\end{equation}

This is equation 1 in Dynare.

Its economic content is the intertemporal substitution mechanism. Define the
expected ex-ante real-rate deviation as

\begin{equation}
r_t^{e}
\equiv
i_t-\mathbb E_t\pi_{t+1}.
\end{equation}

Then

\[
r_t^e\downarrow
\quad\Longrightarrow\quad
c_t-\mathbb E_tc_{t+1}\uparrow.
\]

Thus, if expected inflation rises more than the nominal interest rate, current
consumption becomes relatively attractive. This is the model-consistent
demand-side mechanism behind expenditure being brought forward during rising
inflation.

\subsubsection{Money first-order condition}

Differentiating the Lagrangian with respect to \(M_t\) gives

\begin{equation}
\frac{U_{m,t}}{P_t}
-
\Lambda_t
+
\beta\mathbb E_t\Lambda_{t+1}
=
0.
\end{equation}

Using the bond FOC,

\[
\beta\mathbb E_t\Lambda_{t+1}
=
Q_t^d\Lambda_t,
\]

and equation \eqref{eq:lambda},

\begin{equation}
\frac{U_{m,t}}{U_{C,t}}
=
1-Q_t^d
=
\frac{i_t^{\mathrm{lev}}}{1+i_t^{\mathrm{lev}}}.
\label{eq:money_foc_exact}
\end{equation}

For the separable MIU utility in equation \eqref{eq:utility_level},

\begin{equation}
\zeta_m
\left(\frac{M_t}{P_t}\right)^{-\nu}
C_t^\sigma
=
\frac{i_t^{\mathrm{lev}}}{1+i_t^{\mathrm{lev}}}.
\label{eq:miu_exact}
\end{equation}

The exact MIU specification therefore predicts:

\[
C_t\uparrow
\Rightarrow
\text{money demand rises},
\]

and

\[
i_t^{\mathrm{lev}}\uparrow
\Rightarrow
\text{money demand falls}.
\]

\paragraph{Why the Dynare equation is semi-structural.}

A literal log-linearization of equation \eqref{eq:miu_exact} would impose
cross-equation restrictions tying the money-demand coefficients to
\(\sigma\), \(\nu\), and the steady-state nominal rate. The running model does
not impose those restrictions. Instead it retains the sign and economic
structure implied by MIU while allowing the two elasticities to be estimated
or calibrated independently:

\begin{equation}
\boxed{
m_t
=
\eta_{mc}c_t
-
\eta_{mi}i_t.
}
\label{eq:money_final}
\end{equation}

This is equation 4 in Dynare. The coefficients satisfy

\[
\eta_{mc}>0,
\qquad
\eta_{mi}>0.
\]

This deliberate semi-structural implementation is useful because the Brazilian
money-demand response during high inflation is empirically important and need
not be forced to obey a restrictive preference calibration.

%===============================================================================
\subsection{CES Consumption Allocation}
\label{app:ces}
%===============================================================================

\subsubsection{Consumption aggregator}

Aggregate consumption is a CES composite:

\begin{equation}
C_t
=
\left[
(1-\alpha)^{1/\eta}
C_{H,t}^{\frac{\eta-1}{\eta}}
+
\alpha^{1/\eta}
C_{F,t}^{\frac{\eta-1}{\eta}}
\right]^{\frac{\eta}{\eta-1}}.
\label{eq:ces_level}
\end{equation}

Here \(\alpha\in(0,1)\) governs openness and \(\eta>0\) is the elasticity of
substitution between domestic and foreign goods.

\subsubsection{Expenditure-minimization problem}

For a given \(C_t\), the household chooses \(C_{H,t}\) and \(C_{F,t}\) to
minimize

\begin{equation}
P_{H,t}C_{H,t}+P_{F,t}C_{F,t}
\end{equation}

subject to equation \eqref{eq:ces_level}. Let \(\Xi_t\) denote the multiplier
on the CES constraint. The first-order conditions imply

\begin{equation}
\frac{C_{H,t}}{C_{F,t}}
=
\frac{1-\alpha}{\alpha}
\left(
\frac{P_{H,t}}{P_{F,t}}
\right)^{-\eta}.
\end{equation}

Solving the expenditure-minimization problem yields the demand functions

\begin{equation}
C_{H,t}
=
(1-\alpha)
\left(
\frac{P_{H,t}}{P_t}
\right)^{-\eta}
C_t,
\label{eq:ch_level}
\end{equation}

and

\begin{equation}
C_{F,t}
=
\alpha
\left(
\frac{P_{F,t}}{P_t}
\right)^{-\eta}
C_t.
\label{eq:cf_level}
\end{equation}

The exact CPI dual to the aggregator is

\begin{equation}
P_t
=
\left[
(1-\alpha)P_{H,t}^{1-\eta}
+
\alpha P_{F,t}^{1-\eta}
\right]^{\frac{1}{1-\eta}}.
\label{eq:cpi_level}
\end{equation}

%===============================================================================
\subsection{Terms of Trade and CPI Inflation}
\label{app:tot}
%===============================================================================

Define the terms of trade used in the model as

\begin{equation}
S_t
\equiv
\frac{P_{F,t}}{P_{H,t}},
\qquad
s_t\equiv\log S_t.
\label{eq:s_definition}
\end{equation}

With this convention,

\[
s_t\uparrow
\]

means that imported goods become more expensive relative to Brazilian goods.

A first-order approximation of the CPI around the symmetric relative-price
steady state gives

\begin{equation}
p_t
=
(1-\alpha)p_{H,t}
+
\alpha p_{F,t}.
\end{equation}

Since

\[
s_t=p_{F,t}-p_{H,t},
\]

we obtain

\begin{equation}
p_t-p_{H,t}
=
\alpha s_t,
\label{eq:p_pH}
\end{equation}

and

\begin{equation}
p_{F,t}-p_t
=
(1-\alpha)s_t.
\label{eq:pF_p}
\end{equation}

Substituting these relations into the log-linearized demand equations gives

\begin{equation}
\boxed{
c_{H,t}
=
c_t+\eta\alpha s_t,
}
\label{eq:ch_final}
\end{equation}

and

\begin{equation}
\boxed{
c_{F,t}
=
c_t-\eta(1-\alpha)s_t.
}
\label{eq:cf_final}
\end{equation}

These are equations 2 and 3 in Dynare.

\subsubsection{Law of one price}

Imported final goods satisfy

\begin{equation}
P_{F,t}
=
\mathcal E_tP_{F,t}^{*},
\label{eq:lop}
\end{equation}

where \(\mathcal E_t\) is units of domestic currency per unit of foreign
currency. Therefore

\begin{equation}
s_t
=
e_t+p_{F,t}^{*}-p_{H,t}.
\end{equation}

Taking first differences,

\begin{equation}
\boxed{
s_t-s_{t-1}
=
e_t-e_{t-1}
+
\pi_{F,t}^{*}
-
\pi_{H,t}.
}
\label{eq:tot_final}
\end{equation}

This is equation 9 in Dynare.

\subsubsection{CPI inflation identity}

Taking first differences of equation \eqref{eq:p_pH},

\begin{equation}
\boxed{
\pi_t
=
\pi_{H,t}
+
\alpha(s_t-s_{t-1}).
}
\label{eq:cpi_final}
\end{equation}

This is equation 8 in Dynare.

It is an \emph{identity}, not a Phillips curve. It says that CPI inflation can
rise immediately when imported goods become more expensive, even if
domestic-goods inflation has not yet fully adjusted.

%===============================================================================
\subsection{The Generic Foreign-Price Shock}
\label{app:foreign_price}
%===============================================================================

Foreign imported-goods inflation follows

\begin{equation}
\boxed{
\pi_{F,t}^{*}
=
\rho_F\pi_{F,t-1}^{*}
+
\varepsilon_{F,t}.
}
\label{eq:pif_final}
\end{equation}

This is equation 10 in Dynare.

The model contains no oil quantity, oil production technology, or oil-specific
utility term. The historical oil shocks are interpreted as unusually large
positive realizations of \(\varepsilon_{F,t}\). The advantage is modularity:
the same structural model can later be subjected to other external-price
disturbances without rebuilding the production and household blocks.

%===============================================================================
\subsection{Firms and Imported Intermediate Inputs}
\label{app:firms}
%===============================================================================

\subsubsection{Production technology}

A domestic differentiated producer \(j\) uses labor and imported intermediate
inputs:

\begin{equation}
Y_t(j)
=
A_tN_t(j)^{1-\mu}Z_t(j)^\mu,
\qquad
0<\mu<1.
\label{eq:production_level}
\end{equation}

Aggregating and log-linearizing gives

\begin{equation}
\boxed{
y_t
=
a_t+(1-\mu)n_t+\mu z_t.
}
\label{eq:production_final}
\end{equation}

This is equation 5 in Dynare.

\subsubsection{Conditional cost minimization}

For a required output \(Y_t(j)\), the firm minimizes

\begin{equation}
W_tN_t(j)+P_{Z,t}Z_t(j)
\end{equation}

subject to equation \eqref{eq:production_level}. Let \(\Lambda_t^{MC}\) be the
multiplier on the production constraint:

\begin{align}
\mathcal L_t^F
=&
W_tN_t(j)+P_{Z,t}Z_t(j)
\nonumber\\
&+
\Lambda_t^{MC}
\left[
Y_t(j)-A_tN_t(j)^{1-\mu}Z_t(j)^\mu
\right].
\end{align}

The first-order conditions are

\begin{equation}
W_t
=
\Lambda_t^{MC}
(1-\mu)
\frac{Y_t(j)}{N_t(j)},
\label{eq:foc_N_firm}
\end{equation}

and

\begin{equation}
P_{Z,t}
=
\Lambda_t^{MC}
\mu
\frac{Y_t(j)}{Z_t(j)}.
\label{eq:foc_Z_firm}
\end{equation}

The multiplier \(\Lambda_t^{MC}\) is nominal marginal cost. Solving for the unit
cost gives, up to a constant,

\begin{equation}
MC_t
\propto
\frac{
W_t^{1-\mu}P_{Z,t}^{\mu}
}{A_t}.
\label{eq:unitcost}
\end{equation}

\subsubsection{Imported-input demand}

Equation \eqref{eq:foc_Z_firm} implies

\begin{equation}
P_{Z,t}Z_t
=
\mu MC_tY_t.
\label{eq:input_share}
\end{equation}

Define real marginal cost relative to the domestic output price:

\[
mc_t
\equiv
\log\left(\frac{MC_t}{P_{H,t}}\right)
-
\log\left(\frac{\overline{MC}}{\overline{P_H}}\right).
\]

Assume the foreign price of the imported intermediate moves with the generic
foreign imported-goods price:

\begin{equation}
P_{Z,t}
=
\mathcal E_tP_{Z,t}^{*},
\qquad
P_{Z,t}^{*}=P_{F,t}^{*}
\quad\text{in the baseline}.
\label{eq:pz}
\end{equation}

Hence, to first order,

\begin{equation}
p_{Z,t}-p_{H,t}=s_t.
\end{equation}

Taking logs of equation \eqref{eq:input_share} and removing constants,

\begin{equation}
\boxed{
z_t
=
y_t+mc_t-s_t.
}
\label{eq:z_final}
\end{equation}

This is equation 6 in Dynare.

\subsubsection{Real marginal cost}

Taking logs of equation \eqref{eq:unitcost} relative to \(P_{H,t}\),

\begin{equation}
mc_t
=
(1-\mu)(w_t-p_{H,t})
+
\mu(p_{Z,t}-p_{H,t})
-a_t.
\label{eq:mc_step1}
\end{equation}

Use

\[
w_t-p_{H,t}
=
(w_t-p_t)+(p_t-p_{H,t}),
\]

together with equations \eqref{eq:labor_supply_linear} and \eqref{eq:p_pH}:

\begin{equation}
w_t-p_{H,t}
=
\sigma c_t+\varphi n_t+\alpha s_t.
\end{equation}

Substituting into equation \eqref{eq:mc_step1},

\begin{align}
mc_t
&=
(1-\mu)
\left(
\sigma c_t+\varphi n_t+\alpha s_t
\right)
+
\mu s_t-a_t
\\
&=
(1-\mu)
\left(
\sigma c_t+\varphi n_t
\right)
-a_t
+
\left[
\mu+(1-\mu)\alpha
\right]s_t.
\end{align}

Define

\begin{equation}
\boxed{
\lambda_s
\equiv
\mu+(1-\mu)\alpha.
}
\label{eq:lambda_s}
\end{equation}

Then

\begin{equation}
\boxed{
mc_t
=
(1-\mu)
\left(
\sigma c_t+\varphi n_t
\right)
-a_t
+
\lambda_ss_t.
}
\label{eq:mc_final}
\end{equation}

This is equation 7 in Dynare.

The coefficient \(\lambda_s\) has two components:

\[
\underbrace{\mu}_{\text{direct imported-input cost}}
+
\underbrace{(1-\mu)\alpha}_{\text{CPI--domestic-price wage wedge}}.
\]

Thus the imported-price shock reaches domestic firms without putting the shock
directly into the Phillips curve.

%===============================================================================
\subsection{Calvo Pricing with Backward Indexation}
\label{app:calvo}
%===============================================================================

\subsubsection{Demand for an individual differentiated good}

A standard CES aggregator over domestic differentiated goods implies demand

\begin{equation}
Y_t(j)
=
\left(
\frac{P_{H,t}(j)}{P_{H,t}}
\right)^{-\epsilon}
Y_t,
\qquad
\epsilon>1.
\label{eq:firm_demand}
\end{equation}

A fraction \(1-\theta\) of firms can reoptimize at date \(t\). The remaining
fraction \(\theta\) cannot choose a new price.

\subsubsection{Indexation of non-reset prices}

A firm that cannot reoptimize updates its previous price using past domestic
inflation:

\begin{equation}
P_{H,t}(j)
=
P_{H,t-1}(j)
\Pi_{H,t-1}^{\gamma_I}
\bar\Pi_H^{1-\gamma_I}.
\label{eq:indexation_rule}
\end{equation}

The parameter \(\gamma_I\in[0,1]\) measures backward indexation.

\subsubsection{Reset-price problem}

A firm resetting at date \(t\) chooses \(P_{H,t}^{\#}\) to maximize the expected
discounted stream of profits while its price may remain non-optimized:

\begin{align}
\max_{P_{H,t}^{\#}}
\mathbb E_t
\sum_{k=0}^{\infty}
(\beta\theta)^k
\mathcal Q_{t,t+k}
\Big[
P_{H,t}^{\#}\mathcal I_{t,t+k}
-
MC_{t+k}
\Big]
Y_{t+k|t},
\label{eq:reset_problem}
\end{align}

where \(\mathcal I_{t,t+k}\) cumulates the indexation rule from \(t\) to
\(t+k\), \(\mathcal Q_{t,t+k}\) is the relevant stochastic discounting term,
and \(Y_{t+k|t}\) is demand for the firm's good conditional on its reset price.

The full nonlinear first-order condition equates the reset price to a markup
over a discounted sequence of expected marginal costs. Around the normalized
steady state, the first-order representation can be written in terms of the
``inflation gap''

\begin{equation}
\widetilde\pi_{H,t}
\equiv
\pi_{H,t}-\gamma_I\pi_{H,t-1}.
\end{equation}

The hybrid Phillips curve is

\begin{equation}
\widetilde\pi_{H,t}
=
\beta\mathbb E_t\widetilde\pi_{H,t+1}
+
\kappa mc_t.
\label{eq:nkpc_gap}
\end{equation}

Expanding this equation,

\begin{equation}
\pi_{H,t}-\gamma_I\pi_{H,t-1}
=
\beta\mathbb E_t
\left(
\pi_{H,t+1}-\gamma_I\pi_{H,t}
\right)
+\kappa mc_t.
\end{equation}

Collecting current inflation terms,

\begin{equation}
\boxed{
(1+\beta\gamma_I)\pi_{H,t}
=
\gamma_I\pi_{H,t-1}
+
\beta\mathbb E_t\pi_{H,t+1}
+
\kappa mc_t.
}
\label{eq:nkpc_final}
\end{equation}

This is equation 11 in Dynare.

The implemented slope is

\begin{equation}
\boxed{
\kappa
=
\frac{(1-\theta)(1-\beta\theta)}{\theta}.
}
\label{eq:kappa_final}
\end{equation}

\paragraph{Important implementation note on trend inflation.}

The running model uses a positive steady-state CPI inflation benchmark for the
seigniorage calculation, but the Phillips-curve slope in
equation \eqref{eq:kappa_final} is the standard first-order Calvo coefficient.
A full positive-trend-inflation Calvo derivation would modify the reset-price
recursions and generally change the reduced-form coefficients. That extension
is \emph{not} part of the successfully running baseline. The present appendix
therefore documents the implemented model rather than silently replacing it
with a different positive-trend-inflation pricing block.

%===============================================================================
\subsection{Exports and Foreign Demand}
\label{app:exports}
%===============================================================================

A foreign CES demand system would imply that demand for Brazilian output rises
when Brazilian goods become cheaper relative to foreign goods and when foreign
activity increases. The implemented first-order export equation is the
parsimonious reduced form

\begin{equation}
\boxed{
x_t
=
y_t^{*}+\eta_xs_t.
}
\label{eq:exports_final}
\end{equation}

This is equation 12 in Dynare.

Because \(s_t=P_{F,t}/P_{H,t}\) in logs, \(s_t\uparrow\) means foreign goods are
more expensive relative to Brazilian goods. Hence the sign \(+\eta_xs_t\) is
consistent with an improvement in Brazilian price competitiveness.

%===============================================================================
\subsection{Goods-Market Clearing}
\label{app:goodsmarket}
%===============================================================================

Domestic output is absorbed by household demand for domestic goods, exports,
ordinary government consumption, and public/SOE investment:

\begin{equation}
Y_t
=
C_{H,t}
+
X_t
+
G_t
+
I_t^G.
\label{eq:goods_level}
\end{equation}

Define steady-state shares

\begin{equation}
\omega_{CH}
=
\frac{\overline{C_H}}{\bar Y},
\qquad
\omega_X
=
\frac{\bar X}{\bar Y},
\qquad
\omega_G
=
\frac{\bar G}{\bar Y},
\qquad
\omega_{IG}
=
\frac{\overline{I^G}}{\bar Y}.
\end{equation}

Using the weighted-sum rule in equation \eqref{eq:sharelinearization},

\begin{equation}
\boxed{
y_t
=
\omega_{CH}c_{H,t}
+
\omega_Xx_t
+
\omega_Gg_t
+
\omega_{IG}ig_t.
}
\label{eq:goods_final}
\end{equation}

This is equation 13 in Dynare.

The shares satisfy

\begin{equation}
\boxed{
\omega_{CH}+\omega_X+\omega_G+\omega_{IG}=1.
}
\label{eq:domestic_share_identity}
\end{equation}

Imported consumption and imported intermediate inputs do not appear as uses of
\emph{domestic} output. They enter the external resource balance instead.

%===============================================================================
\subsection{Trade Balance}
\label{app:tradebalance}
%===============================================================================

Measure the trade balance in domestic-good units as

\begin{equation}
TB_t
=
X_t
-
S_tC_{F,t}
-
S_tZ_t.
\label{eq:tb_level}
\end{equation}

Let

\begin{equation}
\omega_{CF}
=
\frac{\overline{C_F}}{\bar Y},
\qquad
\omega_Z
=
\frac{\bar Z}{\bar Y}.
\end{equation}

The baseline normalization imposes

\begin{equation}
\omega_X=\omega_{CF}+\omega_Z,
\label{eq:balanced_trade_normalization}
\end{equation}

so the steady-state goods trade balance is zero.

Because the steady-state trade balance is zero, one must not take its logarithm.
Instead define

\begin{equation}
tb_t
\equiv
\frac{TB_t-\overline{TB}}{\bar Y}.
\end{equation}

First-order approximation of equation \eqref{eq:tb_level} gives

\begin{align}
tb_t
&=
\omega_Xx_t
-
\omega_{CF}(c_{F,t}+s_t)
-
\omega_Z(z_t+s_t).
\end{align}

Thus,

\begin{equation}
\boxed{
tb_t
=
\omega_Xx_t
-
\omega_{CF}(c_{F,t}+s_t)
-
\omega_Z(z_t+s_t).
}
\label{eq:tb_final}
\end{equation}

This is equation 14 in Dynare.

The terms \(c_{F,t}+s_t\) and \(z_t+s_t\) are import \emph{value} deviations:
changes in import expenditure reflect both quantities and relative prices.

%===============================================================================
\subsection{Consolidated Public External Debt}
\label{app:externaldebt}
%===============================================================================

\subsubsection{Economic interpretation}

The model deliberately does not use a household foreign bond. External debt is
instead interpreted as consolidated public-sector foreign indebtedness,
including foreign liabilities associated with the government and relevant
public/SOE channels.

Let the benchmark external-debt-to-output ratio be \(\bar d\). The Dynare
variable \(d_t\) is the deviation of that ratio from its benchmark.

\subsubsection{Motivating debt law}

A simple external debt accumulation relation is

\begin{equation}
D_t^{*}
=
\left(
1+r_{t-1}^{*}+RP_{t-1}
\right)
D_{t-1}^{*}
-
TB_t
+
\text{steady-state normalization terms}.
\label{eq:external_debt_nonlinear_motivation}
\end{equation}

The omitted constant normalization absorbs the steady rollover/resource-transfer
component required to support the chosen positive benchmark \(\bar d\).
Because the model is written entirely in deviations, that constant disappears
from the linear system.

Linearizing the debt ratio around \(\bar d\), and writing the steady world
servicing component as \(\bar r^{*}\), gives the implemented law

\begin{equation}
\boxed{
d_t
=
(1+\bar r^{*})d_{t-1}
-
tb_t
+
\bar d
\left(
i_{t-1}^{*}+rp_{t-1}
\right).
}
\label{eq:d_final}
\end{equation}

This is equation 15 in Dynare.

The signs are central:

\[
tb_t<0
\quad\Rightarrow\quad
d_t\uparrow,
\]

and

\[
i_{t-1}^{*}+rp_{t-1}\uparrow
\quad\Rightarrow\quad
d_t\uparrow.
\]

Thus a trade deficit and higher inherited borrowing costs both increase public
external indebtedness.

\paragraph{Accounting caution.}

Equation \eqref{eq:d_final} is a local debt-ratio law, not a complete
international balance-of-payments system with private capital flows, reserve
accumulation, valuation effects, and GDP-growth denominator effects. Those
objects are intentionally outside the baseline. This is the precise sense in
which the external block is parsimonious.

%===============================================================================
\subsection{Country-Risk Premium}
\label{app:riskpremium}
%===============================================================================

The risk premium is not an optimizing household condition. It is a pricing
wedge imposed on Brazilian external borrowing.

A nonlinear motivation is

\begin{equation}
1+RP_t
=
\exp
\left\{
\phi_D(D_t-\bar D)
+
\phi_{\Delta D}(D_t-D_{t-1})
+
\varepsilon_{rp,t}
\right\}.
\end{equation}

A first-order approximation around zero risk-premium deviation gives

\begin{equation}
\boxed{
rp_t
=
\phi_Dd_t
+
\phi_{\Delta D}(d_t-d_{t-1})
+
\varepsilon_{rp,t}.
}
\label{eq:rp_final}
\end{equation}

This is equation 16 in Dynare.

The first term makes financing conditions sensitive to the level of external
debt. The second is an accelerator: a rapid increase in debt can worsen
financing conditions even for a given debt level.

The model does not assume that this accelerator is mechanically enormous. Its
size is an empirical question.

%===============================================================================
\subsection{World Interest Rate}
\label{app:worldrate}
%===============================================================================

The world nominal interest-rate deviation follows

\begin{equation}
\boxed{
i_t^{*}
=
\rho_{i^{*}}i_{t-1}^{*}
+
\varepsilon_{i^{*},t}.
}
\label{eq:istar_final}
\end{equation}

This is equation 17 in Dynare.

Separating \(i_t^{*}\) from \(rp_t\) is historically important. A deterioration
in Brazil-specific risk and a tightening of global financial conditions are
distinct shocks.

%===============================================================================
\subsection{Sovereign Uncovered Interest Parity}
\label{app:uip}
%===============================================================================

Because households do not trade a foreign bond, UIP is interpreted as an
international-arbitrage condition for Brazilian sovereign/public liabilities.

A foreign investor comparing a one-period domestic-currency Brazilian asset
with a foreign-currency return requires, abstracting from higher-order risk
terms,

\begin{equation}
1+i_t^{\mathrm{lev}}
\simeq
(1+i_t^{*,\mathrm{lev}})
(1+RP_t)
\mathbb E_t
\left(
\frac{\mathcal E_{t+1}}{\mathcal E_t}
\right).
\label{eq:uip_gross}
\end{equation}

Taking logs and linearizing,

\begin{equation}
\boxed{
i_t
=
i_t^{*}
+
rp_t
+
\mathbb E_te_{t+1}
-
e_t.
}
\label{eq:uip_final}
\end{equation}

This is equation 18 in Dynare.

A higher domestic interest rate can therefore reflect:

\begin{enumerate}
    \item higher world rates;
    \item higher Brazilian country risk; or
    \item expected depreciation of the domestic currency.
\end{enumerate}

The equation is called ``sovereign UIP'' in this model precisely to avoid the
incorrect statement that it comes from household portfolio optimization.

%===============================================================================
\subsection{Monetary Policy}
\label{app:monetarypolicy}
%===============================================================================

Monetary policy is specified directly rather than derived from optimization.
The rule is

\begin{equation}
\boxed{
i_t
=
\rho_i i_{t-1}
+
(1-\rho_i)
\left[
\phi_{\pi}\pi_t
+
\phi_E(i_t^{*}+rp_t)
\right]
+
\varepsilon_{i,t}.
}
\label{eq:policy_final}
\end{equation}

This is equation 19 in Dynare.

The final baseline imposes

\begin{equation}
0<\phi_{\pi}<1,
\qquad
\phi_y=0,
\qquad
\phi_E>0.
\end{equation}

Therefore:

\begin{itemize}
    \item the central bank responds only weakly to inflation;
    \item there is no direct output-gap reaction;
    \item external financing conditions matter for the domestic rate.
\end{itemize}

The coefficient \(\phi_{\pi}<1\) is not merely a numerical trick. It is the
model's explicit representation of weak anti-inflation monetary stabilization
during the historical period. Whether the full equilibrium is determinate is
nevertheless a property of the \emph{whole model}, not of \(\phi_\pi\) alone.

%===============================================================================
\subsection{Seigniorage}
\label{app:seigniorage}
%===============================================================================

Nominal money creation is

\[
M_t-M_{t-1}.
\]

Real seigniorage is

\begin{equation}
SG_t
=
\frac{M_t-M_{t-1}}{P_t}.
\label{eq:sg_level_1}
\end{equation}

Define real balances

\[
\mathcal M_t\equiv\frac{M_t}{P_t}
\]

and gross inflation

\[
\Pi_t\equiv\frac{P_t}{P_{t-1}}.
\]

Then

\begin{align}
SG_t
&=
\frac{M_t}{P_t}
-
\frac{M_{t-1}}{P_t}
\\
&=
\mathcal M_t
-
\frac{\mathcal M_{t-1}}{\Pi_t}.
\label{eq:sg_exact}
\end{align}

Let

\[
\omega_M
\equiv
\frac{\overline{\mathcal M}}{\bar Y},
\qquad
\bar\Pi>1.
\]

Write

\[
\mathcal M_t=\overline{\mathcal M}e^{m_t},
\qquad
\Pi_t=\bar\Pi e^{\pi_t}.
\]

Using

\[
e^{m_t}\simeq1+m_t
\]

and

\[
\frac{e^{m_{t-1}}}{\bar\Pi e^{\pi_t}}
\simeq
\frac{1}{\bar\Pi}
\left(
1+m_{t-1}-\pi_t
\right),
\]

subtract the steady-state value and divide by \(\bar Y\). The result is

\begin{equation}
\boxed{
sg_t
=
\omega_M
\left[
m_t
-
\bar\Pi^{-1}m_{t-1}
+
\bar\Pi^{-1}\pi_t
\right].
}
\label{eq:sg_final}
\end{equation}

This is equation 20 in Dynare.

Define

\[
\iota_\Pi\equiv\bar\Pi^{-1}.
\]

Then the exact Dynare implementation is

\[
sg_t
=
\texttt{money\_y}
\left[
m_t
-
\texttt{invPi}\,m_{t-1}
+
\texttt{invPi}\,\pi_t
\right].
\]

Seigniorage is therefore endogenous. It responds to both real money balances
and inflation. The model does not impose an equation such as
\(SG_t=\omega\,Def_t\).

%===============================================================================
\subsection{Taxes, Government Consumption, and Public Investment}
\label{app:fiscal_processes}
%===============================================================================

\subsubsection{Taxes}

Tax revenue as a share of output follows

\begin{equation}
\boxed{
\tau_t
=
\rho_\tau\tau_{t-1}
+
\varepsilon_{\tau,t}.
}
\label{eq:tau_final}
\end{equation}

This is equation 21 in Dynare.

There is no term of the form

\[
\phi_{\tau b}b_{t-1}.
\]

Thus, in the post-1973 baseline,

\begin{equation}
\boxed{
\phi_{\tau b}=0.
}
\end{equation}

Taxes are a persistent revenue source, but debt does not automatically trigger
a stabilizing tax increase.

\subsubsection{Ordinary government consumption}

\begin{equation}
\boxed{
g_t
=
\rho_Gg_{t-1}
+
\varepsilon_{G,t}.
}
\label{eq:g_final}
\end{equation}

This is equation 22 in Dynare.

\subsubsection{Public/SOE investment}

\begin{equation}
\boxed{
ig_t
=
\rho_{IG}ig_{t-1}
+
\varepsilon_{IG,t}.
}
\label{eq:ig_final}
\end{equation}

This is equation 23 in Dynare.

The shock \(\varepsilon_{IG,t}\) is a development-policy shock. It is
deliberately distinct from \(\varepsilon_{F,t}\). The model can therefore ask
what would have happened if Brazil had faced the same foreign-price shock
without the same public-investment response.

%===============================================================================
\subsection{Domestic Public Debt and the Consolidated Financing Identity}
\label{app:domesticdebt}
%===============================================================================

\subsubsection{Economic financing logic}

Domestic debt is the residual financing instrument after accounting for taxes,
seigniorage, and net new external financing.

A useful real-ratio representation is

\begin{equation}
B_t^R
=
(1+r_{d,t-1}^R)B_{t-1}^R
+
G_t
+
I_t^G
-
T_t
-
SG_t
-
NF_t^{*},
\label{eq:domestic_debt_motivation}
\end{equation}

where \(NF_t^{*}\) denotes \emph{net new external financing available for
current fiscal uses}. It is not the same thing as the total increase in
external debt because part of the increase merely capitalizes external
interest.

Write external interest service approximately as

\begin{equation}
INT_t^{*}
\simeq
\bar r^{*}D_{t-1}
+
\bar d
\left(
i_{t-1}^{*}+rp_{t-1}
\right).
\label{eq:external_interest_linear}
\end{equation}

Then net new external financing is

\begin{equation}
NF_t^{*}
=
(D_t-D_{t-1})-INT_t^{*}.
\label{eq:net_external_financing}
\end{equation}

Substituting equation \eqref{eq:net_external_financing} into
\eqref{eq:domestic_debt_motivation} explains why the domestic financing law
contains both

\[
-(d_t-d_{t-1})
\]

and positive external-interest-service terms.

\subsubsection{Real burden of inherited domestic nominal debt}

The real servicing cost of domestic nominal liabilities varies with the
nominal rate and inflation. To first order,

\begin{equation}
r_{d,t-1}^R
\simeq
\bar r_d
+
i_{t-1}
-
\pi_t.
\end{equation}

When this time-varying real rate multiplies the benchmark debt ratio \(\bar b\),
its first-order contribution is

\begin{equation}
\bar b(i_{t-1}-\pi_t).
\end{equation}

Thus current inflation erodes the real value of inherited domestic nominal
liabilities.

\subsubsection{Final domestic-debt equation}

Combining the preceding components and using the steady-state expenditure
shares gives

\begin{equation}
\boxed{
\begin{aligned}
b_t
={}&
(1+\bar r_d)b_{t-1}
+
\omega_Gg_t
+
\omega_{IG}ig_t
-
\tau_t
-
sg_t
\\
&-
(d_t-d_{t-1})
+
\bar r^{*}d_{t-1}
+
\bar d
\left(
i_{t-1}^{*}+rp_{t-1}
\right)
+
\bar b
\left(
i_{t-1}-\pi_t
\right).
\end{aligned}
}
\label{eq:b_final}
\end{equation}

This is equation 24 in Dynare.

The equation can be read term by term:

\begin{align*}
(1+\bar r_d)b_{t-1}
&:
\text{ inherited domestic debt and benchmark servicing},\\
+\omega_Gg_t
&:
\text{ ordinary spending pressure},\\
+\omega_{IG}ig_t
&:
\text{ public-investment pressure},\\
-\tau_t
&:
\text{ tax financing},\\
-sg_t
&:
\text{ monetary financing},\\
-(d_t-d_{t-1})
&:
\text{ gross increase in external financing reduces domestic need},\\
+\bar r^*d_{t-1}
&:
\text{ benchmark external interest is not fresh financing},\\
+\bar d(i_{t-1}^*+rp_{t-1})
&:
\text{ deviations in external borrowing costs},\\
+\bar b(i_{t-1}-\pi_t)
&:
\text{ deviation in the real cost of inherited domestic nominal debt}.
\end{align*}

%===============================================================================
\subsection{Technology and Foreign Output}
\label{app:real_shocks}
%===============================================================================

Productivity follows

\begin{equation}
\boxed{
a_t
=
\rho_aa_{t-1}
+
\varepsilon_{a,t},
}
\label{eq:a_final}
\end{equation}

which is equation 25 in Dynare.

Foreign activity follows

\begin{equation}
\boxed{
y_t^{*}
=
\rho_{y^{*}}y_{t-1}^{*}
+
\varepsilon_{y^{*},t},
}
\label{eq:ystar_final}
\end{equation}

which is equation 26.

%===============================================================================
\subsection{The Exact 26-Equation Dynare System}
\label{app:fullsystem}
%===============================================================================

For reference, the model actually solved by Dynare is the following:

\begin{align}
c_t
&=
\mathbb E_tc_{t+1}
-\frac{1}{\sigma}
\left(
i_t-\mathbb E_t\pi_{t+1}
\right),
\tag{1}
\\
c_{H,t}
&=
c_t+\eta\alpha s_t,
\tag{2}
\\
c_{F,t}
&=
c_t-\eta(1-\alpha)s_t,
\tag{3}
\\
m_t
&=
\eta_{mc}c_t-\eta_{mi}i_t,
\tag{4}
\\
y_t
&=
a_t+(1-\mu)n_t+\mu z_t,
\tag{5}
\\
z_t
&=
y_t+mc_t-s_t,
\tag{6}
\\
mc_t
&=
(1-\mu)(\sigma c_t+\varphi n_t)-a_t+\lambda_ss_t,
\tag{7}
\\
\pi_t
&=
\pi_{H,t}+\alpha(s_t-s_{t-1}),
\tag{8}
\\
s_t-s_{t-1}
&=
e_t-e_{t-1}+\pi_{F,t}^{*}-\pi_{H,t},
\tag{9}
\\
\pi_{F,t}^{*}
&=
\rho_F\pi_{F,t-1}^{*}+\varepsilon_{F,t},
\tag{10}
\\
(1+\beta\gamma_I)\pi_{H,t}
&=
\gamma_I\pi_{H,t-1}
+\beta\mathbb E_t\pi_{H,t+1}
+\kappa mc_t,
\tag{11}
\\
x_t
&=
y_t^{*}+\eta_xs_t,
\tag{12}
\\
y_t
&=
\omega_{CH}c_{H,t}
+\omega_Xx_t
+\omega_Gg_t
+\omega_{IG}ig_t,
\tag{13}
\\
tb_t
&=
\omega_Xx_t
-\omega_{CF}(c_{F,t}+s_t)
-\omega_Z(z_t+s_t),
\tag{14}
\\
d_t
&=
(1+\bar r^{*})d_{t-1}
-tb_t
+\bar d(i_{t-1}^{*}+rp_{t-1}),
\tag{15}
\\
rp_t
&=
\phi_Dd_t+\phi_{\Delta D}(d_t-d_{t-1})
+\varepsilon_{rp,t},
\tag{16}
\\
i_t^{*}
&=
\rho_{i^{*}}i_{t-1}^{*}
+\varepsilon_{i^{*},t},
\tag{17}
\\
i_t
&=
i_t^{*}+rp_t+\mathbb E_te_{t+1}-e_t,
\tag{18}
\\
i_t
&=
\rho_i i_{t-1}
+
(1-\rho_i)
\left[
\phi_\pi\pi_t+\phi_E(i_t^{*}+rp_t)
\right]
+\varepsilon_{i,t},
\tag{19}
\\
sg_t
&=
\omega_M
\left[
m_t-\bar\Pi^{-1}m_{t-1}
+\bar\Pi^{-1}\pi_t
\right],
\tag{20}
\\
\tau_t
&=
\rho_\tau\tau_{t-1}
+\varepsilon_{\tau,t},
\tag{21}
\\
g_t
&=
\rho_Gg_{t-1}
+\varepsilon_{G,t},
\tag{22}
\\
ig_t
&=
\rho_{IG}ig_{t-1}
+\varepsilon_{IG,t},
\tag{23}
\\
b_t
&=
(1+\bar r_d)b_{t-1}
+\omega_Gg_t+\omega_{IG}ig_t-\tau_t-sg_t
\nonumber\\
&\quad
-(d_t-d_{t-1})
+\bar r^{*}d_{t-1}
+\bar d(i_{t-1}^{*}+rp_{t-1})
+\bar b(i_{t-1}-\pi_t),
\tag{24}
\\
a_t
&=
\rho_aa_{t-1}+\varepsilon_{a,t},
\tag{25}
\\
y_t^{*}
&=
\rho_{y^{*}}y_{t-1}^{*}+\varepsilon_{y^{*},t}.
\tag{26}
\end{align}

The 26 endogenous variables are

\begin{equation}
\left\{
c,n,y,z,mc,\pi_H,\pi,s,e,i,m,c_H,c_F,x,tb,d,rp,
i^{*},ig,g,\tau,b,sg,\pi_F^{*},a,y^{*}
\right\}.
\end{equation}

The nine structural shocks are

\begin{equation}
\left\{
\varepsilon_F,
\varepsilon_{i^{*}},
\varepsilon_{rp},
\varepsilon_i,
\varepsilon_{IG},
\varepsilon_G,
\varepsilon_\tau,
\varepsilon_a,
\varepsilon_{y^{*}}
\right\}.
\end{equation}

%===============================================================================
\subsection{Logical Status of Every Equation}
\label{app:equation_types}
%===============================================================================

\begin{description}
\item[Equation 1 -- Consumption Euler equation.] Optimality condition. Derived from the household choice of one-period domestic government bonds.
\item[Equation 2 -- Domestic consumption demand.] Optimal allocation condition. Derived from CES expenditure minimization.
\item[Equation 3 -- Imported consumption demand.] Optimal allocation condition. Derived from CES expenditure minimization.
\item[Equation 4 -- Real money demand.] Semi-structural optimality approximation. Its signs and economic content are motivated by MIU, while its two elasticities are left free in the implemented model.
\item[Equation 5 -- Production technology.] Technology identity. Cobb--Douglas production using labor and imported intermediate inputs.
\item[Equation 6 -- Imported-input demand.] Firm optimality condition. Derived from conditional cost minimization.
\item[Equation 7 -- Real marginal cost.] Derived equilibrium relation. Combines firm unit cost, household labor supply, and relative-price identities.
\item[Equation 8 -- CPI inflation.] Identity. Follows from the CPI and terms-of-trade definitions.
\item[Equation 9 -- Terms of trade.] Identity. Follows from the law of one price and relative-price definitions.
\item[Equation 10 -- Foreign-price process.] Stochastic specification. AR(1) exogenous process.
\item[Equation 11 -- Indexed NKPC.] Firm optimality and price-setting relation. Derived from Calvo reset pricing with backward indexation.
\item[Equation 12 -- Export demand.] Semi-structural foreign-demand equation. Depends on foreign activity and relative-price competitiveness.
\item[Equation 13 -- Goods market.] Market-clearing identity. Equates domestic output with its domestic and foreign uses.
\item[Equation 14 -- Trade balance.] Accounting identity. Export value minus final and intermediate import values.
\item[Equation 15 -- External public debt.] Semi-structural accounting law. Links the trade balance and inherited external financing costs to public external debt.
\item[Equation 16 -- Country-risk premium.] Pricing wedge. Depends on the debt level, debt acceleration, and an autonomous risk shock.
\item[Equation 17 -- World interest rate.] Stochastic specification. AR(1) exogenous process.
\item[Equation 18 -- Sovereign UIP.] Arbitrage condition. Prices domestic and foreign-currency public liabilities.
\item[Equation 19 -- Monetary policy.] Policy rule. Combines a passive inflation response with a response to external financing conditions.
\item[Equation 20 -- Seigniorage.] Accounting identity. Real value of nominal money creation.
\item[Equation 21 -- Taxes.] Fiscal process. Persistent revenue ratio with zero endogenous debt feedback.
\item[Equation 22 -- Government consumption.] Fiscal process. Persistent ordinary government-spending process.
\item[Equation 23 -- Public investment.] Fiscal process. Persistent public/SOE development-policy expenditure.
\item[Equation 24 -- Domestic public debt.] Consolidated financing identity. Domestic debt is residual financing after taxes, seigniorage, and net external financing.
\item[Equation 25 -- Technology.] Stochastic specification. AR(1) productivity process.
\item[Equation 26 -- Foreign output.] Stochastic specification. AR(1) foreign-demand process.
\end{description}

%===============================================================================
\subsection{Baseline Parameterization Used in the Running Model}
\label{app:parameters}
%===============================================================================

The following table records the exact starting values in the running Dynare
file. These are debugging/baseline values rather than claims that all parameters
have already been historically estimated.

\begin{description}
\item[\texttt{beta} = 0.99.] Calibrated quarterly discount factor.
\item[\texttt{sigma} = 1.50.] Calibrated CRRA coefficient / inverse intertemporal elasticity of substitution.
\item[\texttt{varphi} = 1.00.] Calibrated inverse Frisch parameter.
\item[\texttt{alpha} = 0.10.] Calibrated openness/import-share parameter.
\item[\texttt{eta} = 0.80.] Calibrated elasticity between domestic and foreign consumption goods.
\item[\texttt{mu} = 0.12.] Calibrated imported-intermediate production share.
\item[\texttt{lambda\_s} = 0.208.] Derived as \(\mu+(1-\mu)\alpha\); it is not independently estimated.
\item[\texttt{theta} = 0.75.] Baseline value for the Calvo non-reset probability; intended for estimation.
\item[\texttt{gammaI} = 0.70.] Baseline backward-indexation coefficient; intended for estimation.
\item[\texttt{kappa} = 0.0858333.] Derived as \(((1-\theta)(1-\beta\theta))/\theta\).
\item[\texttt{eta\_x} = 0.70.] Calibrated export relative-price elasticity.
\item[\texttt{cH\_y} = 0.68.] Calibrated domestic-household-absorption share of output.
\item[\texttt{x\_y} = 0.12.] Calibrated export share of output.
\item[\texttt{g\_y} = 0.14.] Calibrated ordinary-government-consumption share of output.
\item[\texttt{ig\_y} = 0.06.] Calibrated public/SOE-investment share of output.
\item[\texttt{cf\_y} = 0.05.] Calibrated imported-final-consumption share of output.
\item[\texttt{z\_y} = 0.07.] Calibrated imported-intermediate-expenditure share of output.
\item[\texttt{rhoF} = 0.75.] Baseline foreign-price persistence; intended for estimation.
\item[\texttt{rhoiw} = 0.85.] Baseline world-interest-rate persistence; intended for estimation.
\item[\texttt{phiD} = 0.030.] Baseline debt-level effect on the country-risk premium; intended for estimation.
\item[\texttt{phiDD} = 0.050.] Baseline debt-acceleration effect on the country-risk premium; intended for estimation.
\item[\texttt{dbar} = 0.25.] Calibrated benchmark external public-debt/GDP ratio.
\item[\texttt{rstar\_ss} = 0.0125.] Calibrated quarterly world servicing benchmark.
\item[\texttt{rhoi} = 0.75.] Baseline domestic interest-rate smoothing; intended for estimation.
\item[\texttt{phiPi} = 0.35.] Baseline weak monetary response to inflation; intended for estimation.
\item[\texttt{phiE} = 0.65.] Baseline monetary response to external financing conditions; intended for estimation.
\item[\texttt{eta\_mc} = 0.75.] Baseline activity elasticity of money demand; intended for estimation.
\item[\texttt{eta\_mi} = 4.00.] Baseline interest semi-elasticity of money demand; intended for estimation.
\item[\texttt{money\_y} = 0.08.] Calibrated steady-state real-money-balances/output ratio.
\item[\texttt{Pi\_ss} = \((1.20)^{1/4}\).] Calibrated gross quarterly inflation benchmark corresponding to 20 percent annual inflation.
\item[\texttt{invPi} = \(1/\bar\Pi\).] Derived inverse gross steady-state inflation.
\item[\texttt{rhotau} = 0.85.] Baseline persistence of the effective tax ratio; intended for estimation.
\item[\texttt{rhoG} = 0.80.] Baseline government-consumption persistence; intended for estimation.
\item[\texttt{rhoIG} = 0.95.] Baseline public-investment persistence; intended for estimation.
\item[\texttt{rhoa} = 0.90.] Baseline technology persistence; intended for estimation.
\item[\texttt{rhoystar} = 0.85.] Baseline foreign-output persistence; intended for estimation.
\item[\texttt{rdom\_ss} = 0.010.] Calibrated quarterly domestic-debt servicing benchmark.
\item[\texttt{bbar} = 0.20.] Calibrated benchmark domestic public-debt/GDP ratio.
\end{description}

Two accounting checks hold in the baseline:

\begin{equation}
0.68+0.12+0.14+0.06=1,
\end{equation}

and

\begin{equation}
0.05+0.07=0.12.
\end{equation}

Thus domestic-output absorption shares sum to unity and baseline import
expenditure equals exports.

%===============================================================================
\subsection{Shock Standard Deviations in the Simulation File}
\label{app:shock_sd}
%===============================================================================

The simulation file uses provisional standard deviations solely to generate
one-standard-deviation impulse responses:

\begin{center}
\begin{tabular}{lll}
\hline
Shock & Dynare name & Baseline standard deviation \\
\hline
Foreign price & \(\varepsilon_F\) / \texttt{eps\_pf} & \(0.020\) \\
World interest rate & \(\varepsilon_{i^*}\) / \texttt{eps\_iw} & \(0.010\) \\
Country risk & \(\varepsilon_{rp}\) / \texttt{eps\_rp} & \(0.005\) \\
Domestic monetary policy & \(\varepsilon_i\) / \texttt{eps\_i} & \(0.005\) \\
Public investment & \(\varepsilon_{IG}\) / \texttt{eps\_ig} & \(0.020\) \\
Government consumption & \(\varepsilon_G\) / \texttt{eps\_g} & \(0.010\) \\
Tax revenue & \(\varepsilon_\tau\) / \texttt{eps\_tau} & \(0.005\) \\
Productivity & \(\varepsilon_a\) / \texttt{eps\_a} & \(0.010\) \\
Foreign demand & \(\varepsilon_{y^*}\) / \texttt{eps\_ystar} & \(0.010\) \\
\hline
\end{tabular}
\end{center}

These are not final empirical estimates.

%===============================================================================
\subsection{Steady-State Normalizations}
\label{app:steady_state}
%===============================================================================

Because the model is entered into Dynare using \texttt{model(linear)}, every
endogenous variable in the code has steady state zero. That does \emph{not}
mean the economic steady-state levels are zero.

For example,

\[
d_t=0
\]

in Dynare means

\[
D_t/Y_t=\bar d,
\]

and

\[
b_t=0
\]

means

\[
B_t/Y_t=\bar b.
\]

Similarly, the CPI inflation variable is a deviation around the benchmark gross
inflation rate

\begin{equation}
\bar\Pi=(1.20)^{1/4}.
\end{equation}

Numerically,

\begin{equation}
\bar\Pi \approx 1.0466,
\end{equation}

so the benchmark is about \(4.66\%\) gross quarterly inflation, which
compounds to \(20\%\) annually:

\begin{equation}
\bar\Pi^4=1.20.
\end{equation}

The model's steady-state goods-trade normalization is

\[
\overline{TB}=0.
\]

This does not imply that a country with positive foreign debt literally pays no
interest in the full national accounts. Constant rollover/resource-transfer
terms required to support the benchmark external debt ratio are absorbed into
the local steady-state normalization of the parsimonious external debt law.

%===============================================================================
\subsection{Economic Transmission Mechanisms}
\label{app:transmission}
%===============================================================================

\subsubsection{Foreign-price shock}

A positive shock

\[
\varepsilon_{F,t}>0
\]

raises foreign imported-goods inflation:

\[
\pi_{F,t}^{*}\uparrow.
\]

From the terms-of-trade identity,

\[
s_t\uparrow
\]

unless the exchange rate or domestic inflation offsets the movement
immediately.

There are then three first-round channels.

\paragraph{Channel 1: CPI inflation.}

Equation \eqref{eq:cpi_final} gives

\[
\Delta s_t\uparrow
\Rightarrow
\pi_t\uparrow.
\]

\paragraph{Channel 2: domestic marginal cost.}

Equation \eqref{eq:mc_final} gives

\[
s_t\uparrow
\Rightarrow
mc_t\uparrow
\]

because \(\lambda_s>0\).

The indexed NKPC then implies

\[
mc_t\uparrow
\Rightarrow
\pi_{H,t}\uparrow.
\]

\paragraph{Channel 3: trade balance.}

More expensive imports raise final and intermediate import expenditure:

\[
s_t\uparrow
\Rightarrow
\omega_{CF}(c_{F,t}+s_t)
+
\omega_Z(z_t+s_t)
\uparrow,
\]

which tends to reduce \(tb_t\).

A trade deficit raises external debt:

\[
tb_t\downarrow
\Rightarrow
d_t\uparrow.
\]

Higher debt raises country risk:

\[
d_t\uparrow
\Rightarrow
rp_t\uparrow.
\]

\subsubsection{Inflation persistence}

Even if the original foreign-price shock begins to disappear, equation
\eqref{eq:nkpc_final} contains

\[
\gamma_I\pi_{H,t-1}.
\]

Therefore past inflation remains a determinant of current price setting.

The model thus distinguishes the \emph{impact} mechanism from the
\emph{propagation} mechanism:

\[
\text{foreign price}
\rightarrow
\text{marginal cost}
\rightarrow
\text{initial inflation},
\]

while

\[
\text{indexation}
\rightarrow
\text{persistence}.
\]

\subsubsection{Public-investment response}

A positive

\[
\varepsilon_{IG,t}>0
\]

raises persistent public/SOE investment. Through the goods market it supports
domestic demand. Through the fiscal financing identity it raises the financing
requirement.

Because taxes have no debt feedback, the burden is absorbed through the
combination of seigniorage, external financing, and residual domestic debt.

This makes the counterfactual

\[
\varepsilon_{F,t}>0,
\qquad
\varepsilon_{IG,t}=0
\]

economically meaningful: it represents the same external-price disturbance
without the same development-policy expenditure response.

\subsubsection{Why the second external-price shock can be larger}

By the late 1970s, the economy can enter the second shock with

\[
d_t>0
\]

relative to the initial benchmark and with inflation already elevated through
indexation.

If the second foreign-price shock coincides with

\[
\varepsilon_{i^{*},t}>0,
\]

the external debt-service channel strengthens:

\[
i_t^{*}\uparrow
\Rightarrow
d_{t+1}\uparrow
\Rightarrow
rp_{t+1}\uparrow.
\]

Hence the same structural type of foreign-price shock can have different
effects because the endogenous state of the economy is different.

%===============================================================================
\subsection{Fiscal Dominance, Inflation Taxation, and What the Model Means}
\label{app:fiscal_dominance}
%===============================================================================

The model should not interpret ``fiscal dominance'' and ``inflation taxation''
as mutually exclusive mechanisms.

The fiscal configuration is characterized by

\[
\phi_{\tau b}=0,
\]

so taxes do not automatically rise to stabilize debt, while monetary policy
satisfies

\[
\phi_\pi<1.
\]

At the same time, seigniorage is

\[
sg_t
=
\omega_M
\left[
m_t-\bar\Pi^{-1}m_{t-1}
+\bar\Pi^{-1}\pi_t
\right].
\]

Thus monetary financing is one way in which the fiscal financing requirement
can be accommodated.

The inflation tax and seigniorage are related but not identical concepts.
Seigniorage is the real revenue obtained through money creation. The inflation
tax describes the erosion of the real value of existing nominal money
balances. In the exact identity

\[
SG_t
=
\mathcal M_t-\frac{\mathcal M_{t-1}}{\Pi_t},
\]

both monetary expansion and the inflation-related erosion of inherited real
balances are visible.

%===============================================================================
\subsection{From Structural Equations to a Linear Rational-Expectations System}
\label{app:linear_system}
%===============================================================================

Dynare can represent the model schematically as

\begin{equation}
A_{-1}x_{t-1}
+
A_0x_t
+
A_{+1}\mathbb E_tx_{t+1}
+
B\varepsilon_t
=
0,
\label{eq:lre_system}
\end{equation}

where \(x_t\) collects the 26 endogenous variables and \(\varepsilon_t\)
collects the nine structural shocks.

The presence of \(A_{+1}\) means the model is forward looking. A valid
rational-expectations solution must select a path consistent with both current
equilibrium conditions and expectations of future variables.

Dynare uses a generalized Schur decomposition to study the generalized
eigenvalues of the system.

%===============================================================================
\subsection{Blanchard--Kahn Conditions and Determinacy}
\label{app:bk}
%===============================================================================

The Blanchard--Kahn logic compares:

\begin{enumerate}
    \item the number of unstable generalized eigenvalues; and
    \item the number of non-predetermined, forward-looking variables.
\end{enumerate}

A unique stable rational-expectations solution requires the appropriate
matching between these objects.

Three outcomes are conceptually distinct:

\begin{description}
    \item[Unique equilibrium:] the number of unstable roots matches the number
    of forward-looking degrees of freedom.
    \item[Indeterminacy:] too few unstable roots exist to pin down all jump
    variables, so multiple equilibria may satisfy the linearized conditions.
    \item[No stable equilibrium:] too many unstable roots exist, so no bounded
    path satisfies all equilibrium conditions.
\end{description}

A historically weak monetary response \(\phi_\pi<1\) does not by itself prove
indeterminacy. Determinacy depends on the full fiscal, monetary, debt, and
forward-looking structure.

This is why the running file uses both

\begin{verbatim}
check;
model_diagnostics;
\end{verbatim}

before impulse-response analysis.

%===============================================================================
\subsection{Impulse Responses}
\label{app:irfs}
%===============================================================================

Suppose the solved first-order law of motion is

\begin{equation}
x_t
=
Tx_{t-1}
+
R\varepsilon_t.
\label{eq:solution_law}
\end{equation}

For a one-time unit shock to component \(j\),

\[
\varepsilon_t=e_j,
\]

the impact response is

\begin{equation}
IRF_j(0)=Re_j.
\end{equation}

One period later,

\begin{equation}
IRF_j(1)=TRe_j,
\end{equation}

and generally,

\begin{equation}
\boxed{
IRF_j(h)
=
T^hRe_j.
}
\label{eq:irf_formula}
\end{equation}

Dynare's \texttt{stoch\_simul(order=1,irf=40)} reports responses to a
one-standard-deviation innovation, not necessarily to a unit shock. Therefore
the amplitude of an IRF depends on the shock standard deviation declared in the
\texttt{shocks} block.

%===============================================================================
\subsection{State-Space Representation for Estimation}
\label{app:state_space}
%===============================================================================

Once the model is estimated, its solution can be written as a transition
equation

\begin{equation}
s_t
=
Ts_{t-1}
+
R\varepsilon_t,
\qquad
\varepsilon_t\sim\mathcal N(0,Q),
\label{eq:transition}
\end{equation}

where \(s_t\) is the model state vector.

Observed data \(y_t^{obs}\) are connected to model variables through a
measurement equation

\begin{equation}
y_t^{obs}
=
Zs_t
+
d
+
v_t,
\qquad
v_t\sim\mathcal N(0,H).
\label{eq:measurement}
\end{equation}

The matrices \(Z\) and \(d\) encode data transformations and steady-state
mappings. Measurement error \(v_t\) should only be added when justified by the
data, not simply to force estimation to work.

%===============================================================================
\subsection{Candidate Observables and Measurement Logic}
\label{app:observables}
%===============================================================================

Natural observables for this model include:

\begin{enumerate}
    \item CPI inflation;
    \item domestic-goods or wholesale inflation if available;
    \item real or detrended output;
    \item domestic nominal interest rate;
    \item real monetary balances;
    \item trade balance as a fraction of GDP;
    \item consolidated public external debt as a fraction of GDP;
    \item public/SOE investment;
    \item effective tax revenue as a fraction of GDP;
    \item foreign/import-price inflation;
    \item world interest rates; and
    \item foreign activity.
\end{enumerate}

The measurement equation must respect the units of the structural model.

For example, if observed quarterly inflation is measured as a decimal rate,

\begin{equation}
\pi_t^{obs}
=
\bar\pi+\pi_t.
\end{equation}

If the data are reported in percent,

\begin{equation}
\pi_t^{obs,\%}
=
100(\bar\pi+\pi_t).
\end{equation}

A debt-to-GDP series maps differently because \(d_t\) is already a ratio
deviation:

\begin{equation}
d_t^{obs}
=
\bar d+d_t.
\end{equation}

The trade balance is likewise a GDP-ratio deviation rather than a logged
quantity.

%===============================================================================
\subsection{Data Transformations}
\label{app:data_transformations}
%===============================================================================

A central econometric rule is:

\begin{quote}
Transform the data to match the model; do not reinterpret the model after
seeing the data.
\end{quote}

For a logged real variable such as output, a common measurement relation is

\begin{equation}
100\left[
\log Y_t^{data}
-
\log \bar Y_t^{trend}
\right]
=
100y_t.
\end{equation}

For an annual inflation rate converted to a gross quarterly rate,

\begin{equation}
\Pi_q
=
(1+\pi_{annual})^{1/4}.
\end{equation}

For the baseline annual inflation rate of \(20\%\),

\begin{equation}
\Pi_q=(1.20)^{1/4}.
\end{equation}

Interest-rate conversions must also respect compounding. If \(i_a\) is an
effective annual rate, the corresponding quarterly rate is

\begin{equation}
i_q
=
(1+i_a)^{1/4}-1.
\end{equation}

When historical data are only annual, interpolation does not create genuine
quarterly information. Such series may require mixed-frequency methods,
careful measurement-error treatment, or calibration rather than mechanical
quarterly interpolation.

%===============================================================================
\subsection{Likelihood and the Kalman Filter}
\label{app:kalman}
%===============================================================================

Under Gaussian shocks and a linear state-space representation, the Kalman
filter recursively computes the distribution of the latent state conditional
on observed data.

Let

\[
\widehat s_{t|t-1}
\]

be the predicted state and

\[
P_{t|t-1}
\]

its covariance matrix.

The one-step-ahead forecast error is

\begin{equation}
u_t
=
y_t^{obs}
-
Z\widehat s_{t|t-1}
-d.
\end{equation}

Its covariance is

\begin{equation}
F_t
=
ZP_{t|t-1}Z'
+
H.
\end{equation}

The Gaussian log-likelihood contribution is

\begin{equation}
\ell_t
=
-\frac{1}{2}
\left[
n_y\log(2\pi)
+
\log|F_t|
+
u_t'F_t^{-1}u_t
\right].
\end{equation}

Therefore total log likelihood is

\begin{equation}
\boxed{
\log L(\vartheta)
=
\sum_{t=1}^{T}\ell_t,
}
\label{eq:likelihood}
\end{equation}

where \(\vartheta\) denotes the vector of estimated structural parameters.

%===============================================================================
\subsection{Bayesian Estimation}
\label{app:bayesian}
%===============================================================================

Bayesian estimation combines the likelihood with prior information:

\begin{equation}
\boxed{
p(\vartheta\mid Y)
\propto
L(Y\mid\vartheta)p(\vartheta).
}
\label{eq:bayes}
\end{equation}

Here

\[
p(\vartheta)
\]

is the prior density,

\[
L(Y\mid\vartheta)
\]

is the likelihood generated by the solved DSGE and measurement system, and

\[
p(\vartheta\mid Y)
\]

is the posterior density.

The posterior mode is not the only object of interest. For the key historical
parameters, the paper should report posterior uncertainty.

Examples include

\[
\Pr(\phi_\pi<1\mid Y),
\]

and credible intervals for

\[
\gamma_I,\quad
\phi_D,\quad
\phi_{\Delta D},\quad
\phi_E.
\]

These quantities directly test the mechanisms emphasized by the paper.

%===============================================================================
\subsection{Identification}
\label{app:identification}
%===============================================================================

A parameter is identified when distinct values of that parameter imply
distinguishably different distributions of observables.

The most important identification pairs in this model are:

\begin{description}
    \item[\(\gamma_I\) versus \(\theta\):]
    both affect inflation persistence/dynamics, so inflation alone may not
    sharply separate them.

    \item[\(\phi_D\) versus \(\phi_{\Delta D}\):]
    both affect the risk premium. Debt levels and debt changes must contain
    enough independent variation to identify them separately.

    \item[\(\phi_\pi\) versus \(\rho_i\):]
    interest-rate smoothing can make a policy rule look weak in the short run
    even when its long-run inflation response is larger.

    \item[\(\phi_E\) versus \(\varepsilon_i\):]
    external-financing reactions must be distinguished from autonomous
    domestic policy shocks.

    \item[\(\eta_{mi}\) versus monetary shock volatility:]
    a strong money-demand interest sensitivity and a highly volatile monetary
    shock can both generate large real-balance movements.

    \item[\(\mu\) versus foreign-price shock volatility:]
    imported-input exposure determines transmission strength, while shock
    variance determines the size of the impulse. They should not be allowed to
    substitute freely without data discipline.
\end{description}

Identification should therefore be checked before interpreting posterior
distributions economically.

%===============================================================================
\subsection{Why Some Parameters Are Calibrated and Others Estimated}
\label{app:calibration_estimation}
%===============================================================================

The paper's central mechanism should not be obtained by calibrating the key
parameters to convenient values. Parameters such as

\[
\gamma_I,\quad
\phi_\pi,\quad
\phi_E,\quad
\phi_D,\quad
\phi_{\Delta D}
\]

should, where data permit, face the historical evidence.

By contrast, preference parameters that are weakly identified over a short
historical sample, and accounting shares observed directly in national
accounts, are natural calibration candidates.

Derived quantities such as

\[
\lambda_s=\mu+(1-\mu)\alpha,
\]

\[
\kappa=\frac{(1-\theta)(1-\beta\theta)}{\theta},
\]

and

\[
\bar\Pi^{-1}
\]

must not be estimated independently because doing so would violate the model's
own algebra.

%===============================================================================
\subsection{Forecast-Error Variance Decomposition}
\label{app:fevd}
%===============================================================================

Let the \(h\)-step-ahead forecast error of variable \(j\) be generated by all
structural shocks over the forecast horizon. The forecast-error variance
decomposition measures the fraction of that variance attributable to each
shock.

Conceptually,

\begin{equation}
FEVD_{j,k}(h)
=
\frac{
\text{variance of the }h\text{-step forecast error in }j
\text{ caused by shock }k
}{
\text{total }h\text{-step forecast-error variance of }j
}.
\end{equation}

For this paper, useful questions include:

\begin{itemize}
    \item How much of inflation variance is due to foreign-price shocks?
    \item How much comes from monetary-policy shocks?
    \item How much of external-debt variance is due to trade shocks versus
    world-rate shocks?
    \item Does public-investment policy account for meaningful output and debt
    fluctuations?
\end{itemize}

%===============================================================================
\subsection{Historical Decomposition}
\label{app:historical_decomposition}
%===============================================================================

A historical decomposition reconstructs the observed path of a model variable
as the accumulated contribution of realized structural shocks plus the
contribution of initial conditions.

Schematically,

\begin{equation}
x_t
=
x_t^{init}
+
\sum_k x_{t}^{(\varepsilon_k)}.
\end{equation}

This is especially important for the paper because the historical question is
not merely whether a foreign-price shock \emph{can} raise inflation, but how
much of the actual 1970s inflation path the estimated model attributes to:

\[
\varepsilon_F,\quad
\varepsilon_{IG},\quad
\varepsilon_{i^*},\quad
\varepsilon_{rp},\quad
\varepsilon_i,
\]

and the endogenous propagation of these shocks through indexation and debt.

%===============================================================================
\subsection{Counterfactual Experiments}
\label{app:counterfactuals}
%===============================================================================

The model is particularly well suited to nested counterfactuals.

\subsubsection{No indexation}

Set

\[
\gamma_I=0.
\]

The difference between the baseline and this counterfactual measures how much
persistence is generated by backward price indexation.

\subsubsection{No debt-acceleration premium}

Set

\[
\phi_{\Delta D}=0.
\]

This isolates the incremental contribution of rapidly rising debt to external
financing conditions.

\subsubsection{No imported-input cost channel}

Set

\[
\mu=0.
\]

This removes the direct imported intermediate-input component from marginal
cost, although CPI imported inflation through \(\alpha\) remains.

\subsubsection{No public-investment policy response}

Set the historical public-investment shock to zero:

\[
\varepsilon_{IG,t}=0
\]

during the relevant episode.

\subsubsection{Active monetary stabilization}

For a purely counterfactual experiment, set

\[
\phi_\pi>1.
\]

This asks how the same external shocks would propagate under a much more
aggressive inflation response. It is not the baseline historical regime.

\subsubsection{No world-rate shock in 1979}

Set

\[
\varepsilon_{i^*,t}=0
\]

during the second episode. The difference measures the role of deteriorating
global financing conditions conditional on the inherited debt stock.

%===============================================================================
\subsection{Common Conceptual Errors to Avoid}
\label{app:errors}
%===============================================================================

\begin{enumerate}
    \item \textbf{Do not reintroduce a household foreign bond.}
    External borrowing is consolidated public-sector debt in the final model.

    \item \textbf{Do not call \(s_t\) an oil price.}
    It is the terms of trade \(P_F/P_H\). Oil shocks are represented through
    \(\varepsilon_F\).

    \item \textbf{Do not put \(\varepsilon_F\) directly into the NKPC.}
    It already affects CPI inflation and marginal cost through \(s_t\).

    \item \textbf{Do not add an output gap to the monetary rule.}
    The final historical rule has \(\phi_y=0\).

    \item \textbf{Do not add debt feedback to taxes.}
    The final post-1973 tax process has \(\phi_{\tau b}=0\).

    \item \textbf{Do not treat the trade balance as government revenue.}
    It belongs to the external resource/debt block.

    \item \textbf{Do not log-linearize a zero steady-state trade balance.}
    Use a level deviation relative to steady-state GDP.

    \item \textbf{Do not estimate derived parameters independently.}
    In particular, \(\lambda_s\), \(\kappa\), and \(\bar\Pi^{-1}\) are
    algebraically determined.

    \item \textbf{Do not interpret \(\phi_\pi<1\) as sufficient evidence of
    indeterminacy.}
    Determinacy is a property of the complete rational-expectations system.

    \item \textbf{Do not interpret every equation as a deep structural FOC.}
    The money-demand equation, risk-premium equation, and debt laws contain
    deliberate semi-structural elements.

    \item \textbf{Do not confuse a one-standard-deviation IRF with a unit shock.}
    Dynare scales IRFs by the standard deviation in the \texttt{shocks} block.

    \item \textbf{Do not interpret a large simulated debt path as proof of
    historical accuracy.}
    Historical validation requires estimation, decomposition, and
    counterfactual comparison.
\end{enumerate}

%===============================================================================
\subsection{Dynare-to-Mathematics Dictionary}
\label{app:dynare_dictionary}
%===============================================================================

\begin{center}
\begin{tabular}{lll}
\hline
Dynare & Mathematics & Meaning \\
\hline
\texttt{c} & \(c_t\) & aggregate consumption \\
\texttt{n} & \(n_t\) & labor \\
\texttt{y} & \(y_t\) & domestic output \\
\texttt{z} & \(z_t\) & imported intermediate-input volume \\
\texttt{mc} & \(mc_t\) & real marginal cost \\
\texttt{piH} & \(\pi_{H,t}\) & domestic-goods inflation \\
\texttt{pi} & \(\pi_t\) & CPI inflation \\
\texttt{s} & \(s_t\) & terms of trade, \(\log(P_F/P_H)\) \\
\texttt{e} & \(e_t\) & nominal exchange rate \\
\texttt{i} & \(i_t\) & domestic nominal rate \\
\texttt{m} & \(m_t\) & real money balances \\
\texttt{ch} & \(c_{H,t}\) & domestic household consumption \\
\texttt{cf} & \(c_{F,t}\) & imported household consumption \\
\texttt{x} & \(x_t\) & exports \\
\texttt{tb} & \(tb_t\) & trade balance / GDP deviation \\
\texttt{d} & \(d_t\) & external public debt / GDP deviation \\
\texttt{rp} & \(rp_t\) & country-risk premium \\
\texttt{istar} & \(i_t^*\) & world interest rate \\
\texttt{ig} & \(ig_t\) & public/SOE investment \\
\texttt{g} & \(g_t\) & ordinary government consumption \\
\texttt{tau} & \(\tau_t\) & tax revenue / GDP deviation \\
\texttt{b} & \(b_t\) & domestic public debt / GDP deviation \\
\texttt{sg} & \(sg_t\) & seigniorage / GDP deviation \\
\texttt{pif} & \(\pi_{F,t}^*\) & foreign imported-goods inflation \\
\texttt{a} & \(a_t\) & productivity \\
\texttt{ystar} & \(y_t^*\) & foreign activity \\
\hline
\end{tabular}
\end{center}

%===============================================================================
\subsection{A Step-by-Step Reading Strategy}
\label{app:reading_strategy}
%===============================================================================

When returning to the model after time away, the most efficient order is:

\begin{enumerate}
    \item Start with equations (2), (3), (8), and (9) to reconstruct the
    Gal\'i--Monacelli open-economy price system.

    \item Read equations (5)--(7) to understand how imported intermediate
    inputs convert the foreign relative-price movement into domestic marginal
    cost.

    \item Read equation (11) to understand why the initial cost increase becomes
    persistent domestic inflation.

    \item Read equations (12)--(16) to follow the trade-balance, external-debt,
    and risk-premium feedback.

    \item Read equations (18)--(19) jointly: sovereign UIP explains market
    pricing, while the policy rule describes the domestic authority's behavior.

    \item Read equations (4) and (20) jointly to understand real money demand and
    seigniorage.

    \item Read equations (21)--(24) to understand why taxes do not close the
    fiscal system and why domestic debt becomes the residual financing
    instrument.

    \item Finally, return to equation (1). The Euler equation shows how the
    nominal-rate/inflation combination feeds back into aggregate demand.
\end{enumerate}

%===============================================================================
\subsection{Compact Causal Map}
\label{app:causal_map}
%===============================================================================

The first historical external-price episode is summarized by

\begin{equation}
\varepsilon_{F,t}>0
\Rightarrow
\pi_{F,t}^{*}\uparrow
\Rightarrow
s_t\uparrow.
\end{equation}

Then

\begin{equation}
s_t\uparrow
\Rightarrow
\begin{cases}
\pi_t\uparrow,\\
mc_t\uparrow,\\
tb_t\downarrow.
\end{cases}
\end{equation}

The cost channel gives

\begin{equation}
mc_t\uparrow
\Rightarrow
\pi_{H,t}\uparrow,
\end{equation}

while indexation gives

\begin{equation}
\pi_{H,t}\uparrow
\Rightarrow
\pi_{H,t+1},\pi_{H,t+2},\ldots
\text{ remain elevated}.
\end{equation}

The external channel gives

\begin{equation}
tb_t\downarrow
\Rightarrow
d_t\uparrow
\Rightarrow
rp_t\uparrow.
\end{equation}

The fiscal-development response can simultaneously generate

\begin{equation}
\varepsilon_{IG,t}>0
\Rightarrow
ig_t\uparrow
\Rightarrow
\text{financing requirement}\uparrow.
\end{equation}

Because taxes do not respond to debt and monetary policy reacts weakly to
inflation, adjustment occurs through the interaction of

\[
sg_t,\qquad b_t,\qquad d_t,
\]

with domestic interest rates also influenced by

\[
i_t^*+rp_t.
\]

In the later episode, a positive world-rate shock adds

\begin{equation}
\varepsilon_{i^*,t}>0
\Rightarrow
i_t^*\uparrow
\Rightarrow
\text{external debt service}\uparrow,
\end{equation}

when inherited external debt is already larger.

%===============================================================================
\subsection{Final Consistency Checklist}
\label{app:consistency_checklist}
%===============================================================================

The final baseline is internally defined by the following restrictions:

\begin{align}
&\text{Endogenous variables}=26,\\
&\text{Equations}=26,\\
&\text{Household foreign bond}=0,\\
&\text{Explicit oil variable}=0,\\
&\text{Direct oil/foreign-price shock in NKPC}=0,\\
&\phi_y=0,\\
&\phi_{\tau b}=0,\\
&\phi_\pi<1\ \text{in the historical baseline},\\
&\lambda_s=\mu+(1-\mu)\alpha,\\
&\kappa=\frac{(1-\theta)(1-\beta\theta)}{\theta},\\
&\omega_{CH}+\omega_X+\omega_G+\omega_{IG}=1,\\
&\omega_{CF}+\omega_Z=\omega_X.
\end{align}

Any future extension that changes one of these restrictions should be treated
as a new model version rather than silently folded into the baseline.

%===============================================================================
\subsection{Summary}
\label{app:summary}
%===============================================================================

The model's economic architecture can be summarized in six blocks:

\begin{enumerate}
    \item \textbf{Households:} consumption smoothing, labor supply, domestic
    bond demand, and money demand.

    \item \textbf{Open-economy prices:} a Gal\'i--Monacelli consumption
    aggregator, the terms of trade, and imported CPI inflation.

    \item \textbf{Production and pricing:} imported intermediate inputs raise
    marginal cost when foreign goods become expensive, while indexed Calvo
    pricing propagates domestic inflation.

    \item \textbf{External finance:} trade deficits accumulate consolidated
    public external debt, which feeds into country risk, while world interest
    rates provide a separate global financing shock.

    \item \textbf{Monetary and fiscal policy:} monetary policy responds weakly
    to inflation and to external financing pressure; taxes do not
    automatically stabilize debt; public investment is a separate policy
    process; seigniorage is endogenous.

    \item \textbf{Debt closure:} domestic public debt is the residual financing
    instrument after taxes, seigniorage, and net external financing.
\end{enumerate}

Mathematically, the model is a 26-equation linear rational-expectations system.
Statistically, once measurement equations are added, it becomes a linear
state-space model that can be estimated with likelihood-based or Bayesian
methods and analyzed through impulse responses, forecast-error variance
decompositions, historical decompositions, and counterfactual simulations.

The central empirical question is therefore not whether the model can be made
to generate high inflation. The central question is whether Brazilian data
assign quantitatively important roles to the structural mechanisms represented
by

\[
\gamma_I,\quad
\phi_\pi,\quad
\phi_E,\quad
\phi_D,\quad
\phi_{\Delta D},\quad
\mu,\quad
\rho_{IG},\quad
\eta_{mi},
\]

and whether the full model explains the historical interaction among foreign
prices, inflation inertia, external indebtedness, seigniorage, public
investment, and weak monetary stabilization better than restricted versions
that remove those channels.
